\section{Conclusão: A Bioética da Era Preditiva}
\label{sec:conclusao}

A matriz ético-legal que norteará os gêmeos digitais odontológicos encontra-se ainda em sua gênese. O Brasil ocupa uma posição estratégica: a posse de um sistema universal de saúde de dimensões continentais (SUS), a consolidação da Lei Geral de Proteção de Dados (LGPD) e a vanguarda normativa estabelecida pela ABNT NBR ISO/IEC 3173 \cite{abinc2026brasil} criam as pré-condições para a estruturação de um modelo ético exemplar em nível global.

Entretanto, este arcabouço carece de maturação regulatória para a odontologia preditiva. O desenvolvimento ético não deve ser percebido como um fator restritivo ou burocrático, mas como a única infraestrutura capaz de garantir a aceitação clínica e a confiança pública (\textit{Trust}). Sem transparência algorítmica, sem clareza sobre a responsabilidade civil (\textit{malpractice}), sem o combate ativo aos vieses de dados e sem políticas de equidade distributiva, os gêmeos digitais correm o risco de se tornarem ferramentas formidáveis a serviço exclusivo de uma elite socioeconômica, exacerbando o déficit de saúde bucal e corrompendo o compromisso bioético central da odontologia: o bem-estar do paciente.

O momento demanda o amadurecimento de parcerias institucionais. As decisões regulatórias arquitetadas na presente década delinearão se os gêmeos digitais odontológicos atuarão como agentes de emancipação clínica e precisão ou como caixas-pretas concentradoras de poder tecnológico. A escolha é, fundamentalmente, política, social e inadiável.
